# 问题2：血糖预测GAM模型——完整数学公式

## 一、分布假设

\[
Y_i \mid \boldsymbol{X}_i \sim N(\mu_i, \sigma^2)
\]

其中 \(Y_i\) 为第 \(i\) 个个体的血糖值。

## 二、最终模型方程（交互GAM）

\[
\boxed{Y_i = \beta_0 + f_1(\text{Age}_i) + f_2(\text{TG}_i) + f_{12}(\text{Age}_i, \text{TG}_i) + \sum_{m=3}^{7} f_m(X_{mi}) + \varepsilon_i}
\]

展开为完整形式：

\[
\begin{aligned}
Y_i = \beta_0 &+ f_1(\text{年龄}_i) + f_2(\text{TG}_i) \\
&+ f_{12}(\text{年龄}_i, \text{TG}_i) \\
&+ f_3(\text{RBC}_i) + f_4(\text{MCHC}_i) + f_5(\text{HGB}_i) \\
&+ f_6(\text{ALT}_i) + f_7(\text{性别}_i) + \varepsilon_i
\end{aligned}
\]

## 三、基函数展开（B样条）

每个平滑项用三次B样条基展开：

\[
f_j(X_j) = \sum_{k=1}^{K_j} \beta_{jk} B_{jk}(X_j), \quad K_j = 10
\]

\[
f_{12}(\text{Age}, \text{TG}) = \sum_{k=1}^{K_1}\sum_{l=1}^{K_2} \beta_{kl} B_{1k}(\text{Age}) B_{2l}(\text{TG})
\]

## 四、惩罚机制

\[
L = \sum_{i=1}^{n} \left(Y_i - \sum_{j} f_j(X_{ji})\right)^2 + \sum_{j=1}^{7} \lambda_j \int [f_j''(x)]^2 dx + \lambda_{12} \iint \left(\frac{\partial^2 f_{12}}{\partial \text{Age} \cdot \partial \text{TG}}\right)^2 d\text{Age} \cdot d\text{TG}
\]

- 第一项：残差平方和（RSS）
- 第二项：各主效应的粗糙度惩罚
- 第三项：交互项的张量积惩罚

## 五、约束条件（可识别性）

\[
E[f_j(X_j)] = 0, \quad \forall j = 1, 2, \ldots, 7
\]

\[
E[f_{12}(\text{Age}, \text{TG})] = 0
\]

## 六、P-IRLS求解算法

因采用恒等链接 \(g(\mu_i) = \mu_i\)，每次迭代中：

\[
\tilde{Y} = \hat{\eta}^{(old)} + (Y - \hat{\mu}^{(old)})
\]

权重恒为 \(w = 1\)。求解带罚函数的加权最小二乘：

\[
\hat{\boldsymbol{\beta}}^{(new)} = (\boldsymbol{X}^\top \boldsymbol{W} \boldsymbol{X} + \lambda \boldsymbol{S})^{-1} \boldsymbol{X}^\top \boldsymbol{W} \tilde{Y}
\]

迭代至对数似然变化量 \(< 10^{-7}\) 时判定收敛。

## 七、平滑参数选择

使用REML（受限最大似然）：

\[
\hat{\lambda} = \arg\min_{\lambda} \left\{ -2 \ell(\hat{\boldsymbol{\beta}}(\lambda)) + \log |\boldsymbol{X}^\top \boldsymbol{W} \boldsymbol{X} + \lambda \boldsymbol{S}| \right\}
\]

## 八、模型性能

\[
\begin{aligned}
R^2 &= 0.1618 \\
\text{调整 } R^2 &= 0.1608 \\
\text{AIC} &= 20819.37 \\
\text{RMSE} &= 1.3985 \\
\text{CV-RMSE} &= 1.4264 \pm 0.1330
\end{aligned}
\]

## 九、Gamma GAM升级方案（当异方差显著时）

\[
Y_i \mid \boldsymbol{X}_i \sim \text{Gamma}(\mu_i, \phi)
\]

\[
\text{Var}(Y_i \mid \boldsymbol{X}_i) = \phi \cdot \mu_i^2
\]

\[
\log(\mu_i) = \beta_0 + \sum_{j=1}^{7} f_j(X_{ji}) + f_{12}(\text{Age}_i, \text{TG}_i)
\]
